\chapter{Polyuria}

\section*{Definition}
Polyuria is the excessive production of urine, defined as a urinary output exceeding 3,000 mL per 24 hours in adults or >40 mL/kg/day. It must be distinguished from urinary frequency, which refers to increased voiding episodes without an increase in total volume.

\section*{What Happens in Your Body}
Polyuria arises from either a water diuresis or an osmotic diuresis. Water diuresis results from inadequate secretion of or renal resistance to antidiuretic hormone (ADH), impairing collecting duct water reabsorption. Osmotic diuresis occurs when osmotically active solutes (glucose, urea, mannitol) in tubular fluid generate an osmotic gradient that retains water within the lumen, overwhelming the kidney's concentrating capacity.

\section*{Causes}
\begin{itemize}
  \item Water diuresis: Central diabetes insipidus (ADH deficiency from pituitary/hypothalamic lesions, head trauma, surgery), Nephrogenic diabetes insipidus (lithium toxicity, hypercalcemia, hypokalemia, inherited mutations, chronic kidney disease), Primary polydipsia (psychogenic, dipsogenic)
  \item Osmotic diuresis: Diabetes mellitus (uncontrolled hyperglycemia), Post-obstructive diuresis, Mannitol therapy, High-protein tube feeding (urea-driven)
  \item Pharmacologic causes: Diuretics, SGLT2 inhibitors, Caffeine, Alcohol
  \item Other causes: Hypercalcemia, Hypokalemia, Sickle cell trait/disease, Amyloidosis, Sj\"ogren syndrome
\end{itemize}

\section*{When to See a Doctor}
See your doctor if:
\begin{itemize}
  \item Urine output consistently exceeds 3 L/day
  \item Excessive thirst accompanied by unintended weight loss or fatigue (suspect diabetes mellitus)
  \item Nocturia interfering with sleep or causing dehydration
  \item Post-head trauma or neurosurgical onset of polyuria and polydipsia
  \item Concurrent symptoms: vision changes, headache, or focal neurologic deficits
\end{itemize}

\section*{Related Symptoms}
\begin{itemize}
  \item Polydipsia
  \item Nocturia
  \item Urinary frequency
  \item Dehydration
  \item Weight loss
  \item Fatigue
  \item Dry mouth
  \item Dizziness
\end{itemize}
\textbf{Important:} A 24-hour urine collection with simultaneous serum and urine osmolality, followed by a water deprivation test or desmopressin challenge, distinguishes diabetes insipidus subtypes from primary polydipsia.

\section*{Self-Care / Home Management}
\begin{itemize}
  \item Maintain adequate fluid intake to match urine losses and prevent hypernatremia.
  \item Monitor blood glucose at home if diabetes mellitus is known or suspected.
  \item Avoid excessive caffeine, alcohol, and diuretic substances that worsen polyuria.
  \item Track daily urine output and weight changes to assess volume status.
  \item Seek evaluation if polyuria persists despite normalizing blood glucose.
\end{itemize}

\section*{How Your Doctor Will Evaluate You}
\begin{itemize}
  \item Confirm polyuria with 24-hour urine collection ($>$3 L).
  \item Check serum glucose, electrolytes, calcium, and renal function.
  \item Measure paired serum and urine osmolality to classify as water versus osmotic diuresis.
  \item Perform water deprivation test with desmopressin challenge when diabetes insipidus is suspected.
  \item Brain MRI for suspected central diabetes insipidus; review medications for nephrogenic causes.
\end{itemize}

\section*{Treatment Options}
\begin{itemize}
  \item Central diabetes insipidus: Desmopressin (DDAVP) intranasal, oral, or parenteral.
  \item Nephrogenic diabetes insipidus: Discontinue offending drugs, correct hypercalcemia/hypokalemia; thiazide diuretics with sodium restriction or amiloride may reduce urine output.
  \item Diabetes mellitus: Optimize glycemic control with insulin or oral hypoglycemics.
  \item Primary polydipsia: Behavioral modification, psychiatric support, and fluid restriction.
\end{itemize}

\section*{References}
\begin{thebibliography}{99}
\bibitem{polyuria:polyuria1} Robertson GL. "Diabetes insipidus." N Engl J Med. 2016;375(14):1387-1396.
\bibitem{polyuria:polyuria2} Verbalis JG. "Disorders of body water homeostasis." In: Harrison\'s Principles of Internal Medicine. 21st ed. McGraw-Hill; 2022.
\bibitem{polyuria:polyuria3} Gardella A, Bhatt P. "Polydipsia and polyuria in children." Pediatr Rev. 2018;39(5):251-260.
\bibitem{polyuria:polyuria4} American Diabetes Association. "Classification and diagnosis of diabetes." Diabetes Care. 2022;45(Suppl 1):S17-S38.
\bibitem{polyuria:polyuria5} Gubbi S, Hannah-Shmouni F, Koch CA. "Polyuria." In: StatPearls. StatPearls Publishing; 2023.
\bibitem{polyuria:polyuria6} Berl T, Verbalis JG. "Pathophysiology of disorders of water metabolism." In: Brenner \& Rector\'s The Kidney. 11th ed. Elsevier; 2020.
\end{thebibliography}
